
\subsubsection{5.1 Advantage Variance Comparison (H-M1)}

Table 1 presents the primary comparison between function-level and repository-level GRPO advantage variance across 120 training steps per condition.

\textbf{Table 1.} GRPO advantage variance comparison (CodeLlama-7b-Instruct-hf, G=8, B=4, 120 steps per condition).

\begin{table}[htbp]
\centering
\begin{tabular}{llll}
\toprule
Condition & Mean Adv. Variance & Positive Rate & Steps \\
\midrule
Function-level (APPS+CodeContests) & 0.004167 & ≈0\% & 120 \\
Repo-level (SWE-bench Verified) & 0.316667 & ≈6\% & 120 \\
Ratio (repo / function-level) & 76× & — & — \\
\bottomrule
\end{tabular}
\end{table}

\textbf{Statistical test:} Welch's t-test on log-transformed variance: t=20.37, p=5.34×10⁻⁴⁴, Cohen's d=1.904.

The function-level condition produces advantage variance of 0.004167, which is two orders of magnitude lower than the repository-level condition (0.316667). The ratio is 76×. The statistical test rejects the null hypothesis with t=20.37 and p=5.34×10⁻⁴⁴. Cohen's d=1.904 indicates a very large effect by conventional thresholds (d>0.8).

These results satisfy the gate condition for H-M1 (MUST_WORK gate: PASS).

\begin{figure}[htbp]
\centering
\includegraphics[width=\columnwidth]{figures/advantage_variance_over_steps.png}
\caption{Per-step advantage variance over 120 training steps}
\label{fig:advantage_variance_over_steps}
\end{figure}

\textit{Figure 1. Per-step GRPO advantage variance over 120 training steps for function-level (APPS+CodeContests) and repository-level (SWE-bench Verified) conditions.}

\begin{figure}[htbp]
\centering
\includegraphics[width=\columnwidth]{figures/advantage_variance_bar.png}
\caption{Mean advantage variance bar chart}
\label{fig:advantage_variance_bar}
\end{figure}

\textit{Figure 2. Mean GRPO advantage variance with 95\% CI for both conditions.}

\subsubsection{5.2 Reward Sparsity as Mechanistic Pathway (RQ2)}

The positive reward rate directly explains the variance difference. In the function-level condition, approximately 0\% of completions receive non-zero reward: CodeLlama-7b-Instruct essentially never generates a correct solution to an APPS or CodeContests problem within the allotted computation. Every group therefore has std(r)=0, every advantage is zero, and every gradient contribution vanishes.

In the repository-level condition, approximately 6\% of completions receive non-zero reward through partial credit file-path overlap. This is sufficient to produce non-degenerate advantage groups in a meaningful fraction of training steps.

\begin{figure}[htbp]
\centering
\includegraphics[width=\columnwidth]{figures/positive_rate_over_steps.png}
\caption{Per-step positive reward rate}
\label{fig:positive_rate_over_steps}
\end{figure}

\textit{Figure 3. Per-step positive reward rate (fraction of completions with non-zero reward) over 120 training steps.}

\begin{figure}[htbp]
\centering
\includegraphics[width=\columnwidth]{figures/reward_distribution.png}
\caption{Reward distribution comparison}
\label{fig:reward_distribution}
\end{figure}

\textit{Figure 4. Reward distributions for function-level (left) and repository-level (right) conditions.}

\subsubsection{5.3 Infrastructure Validation (H-E1)}

All four curriculum GRPO conditions completed the 10-step smoke test without errors:

\textbf{Table 2.} H-E1 smoke test results (10 steps per condition).

\begin{table}[htbp]
\centering
\begin{tabular}{llll}
\toprule
Condition & Exit Code & Checkpoints Saved & Reward Density Logs Produced \\
\midrule
curriculum & 0 & Yes & Yes \\
uniform & 0 & Yes & Yes \\
easy_only & 0 & Yes & Yes \\
hard_only & 0 & Yes & Yes \\
\bottomrule
\end{tabular}
\end{table}

Reward density values were 0.0 across all 10 steps for all four conditions. This is consistent with the H-M1 finding: at initialization, even easy-tier APPS problems yield zero reward for the base model. The infrastructure is confirmed to be functioning correctly in that the logging system accurately records these zero values rather than producing spurious non-zero outputs.

The gate condition for H-E1 (MUST_WORK gate: infrastructure passes smoke test) is satisfied. The full 5000-step training run was not executed; no pass@1 comparison between curriculum and uniform conditions is available.

\subsubsection{5.4 H-M2 Analysis (Not Completed)}

H-M2 planned to measure Pearson correlation between per-checkpoint reward density and subsequent pass@1 gain (target: r>0.5) and to compare reward entropy between curriculum and uniform conditions. This analysis requires H-E1 full training logs (4 conditions × 5000 steps). Because full training was not executed, only the curriculum condition has checkpoints (8 checkpoints, steps 500–4000); uniform, easy_only, and hard_only conditions have no checkpoints. All reward density values are 0.0 across all checkpoints. The analysis pipeline was verified correct but could not be executed with the available data. H-M2 gate result: FAILED (SHOULD_WORK gate; not a blocking failure).
